\documentclass[10pt,a4paper]{article}
\usepackage[a4paper,margin=1.3cm,top=1.0cm,bottom=1.1cm]{geometry}
\usepackage[T1]{fontenc}
\usepackage{mathpazo}
\usepackage[scaled=0.92]{helvet}
\usepackage{microtype}
\usepackage{booktabs,tabularx,array}
\usepackage{xcolor}
\usepackage{titlesec}
\usepackage{enumitem}
\usepackage[hidelinks]{hyperref}
\usepackage{caption}

\titleformat{\section}{\sffamily\bfseries\normalsize\color{teal!55!black}}{\thesection}{0.5em}{}
\titlespacing*{\section}{0pt}{5pt}{2pt}
\setlength{\parindent}{0pt}
\setlength{\parskip}{2pt}
\setlist{nosep,leftmargin=1.1em,topsep=1pt}
\captionsetup{font=footnotesize,labelfont=bf,skip=2pt}
\setlength{\textfloatsep}{6pt}
\setlength{\floatsep}{5pt}
\setlength{\intextsep}{5pt}
\setlength{\tabcolsep}{4pt}
\pagestyle{empty}

\begin{document}
\small

{\sffamily\bfseries\Large Lumbar spine segmentation: training on VerSe instead of TotalSegmentator}

\section{Setup}
The pipeline, network (residual-encoder U-Net), loss, optimiser, 11\,000 training steps, checkpoint rule
and test sets are those of the original project. Only the training data change, and with them the voxel size:
\begin{itemize}
\item \textbf{TotalSegmentator-trained} (original project): 902 TotalSegmentator v3 training scans, 1.5\,mm
  voxels (the native TotalSegmentator size), batch 3.
\item \textbf{VerSe-trained} (this project): the VerSe19 training set, 80 scans of 67 patients; 62 usable after
  excluding 15 scans without a lumbar vertebra and 3 with an L6 (the network has no L6 class); 56 for training
  and 6 for choosing the checkpoint, split by patient. Voxel size \textbf{1\,mm}: 39 of the 80 scans are
  already 1\,mm isotropic, whereas at 1.5 or 2\,mm all 80 would be resampled. One 1\,mm patch
  (192$\times$192$\times$320 voxels, the same 19.2$\times$19.2$\times$32\,cm) fills the 24\,GB GPU, so batch 1.
\end{itemize}

% Generated by lumbar/scripts/10_write_short_comparison_report.py -- do not edit by hand.

\begin{table}[h]
\centering\footnotesize
\caption{Anomalies found by the automatic check (Table~1 of the original report, same rules), with the detection and handling of each. Counts are numbers of scans: TS = TotalSegmentator (1939), VS = VerSe19 validation + test (80), \textbf{VS train = VerSe19 training set (80), used here for training}. Upper block: header; lower block: image and label content. $^{a}$Five VS-train scans have exactly 3$\times$ (1\,$\times$\,1\,$\times$\,3\,mm) and are not counted. $^{b}$The network has no L6 class; the 3 VS-train scans with L6 are not used for training.}
\label{tab:anom}
\begin{tabularx}{\linewidth}{@{}l>{\raggedright\arraybackslash}X>{\raggedright\arraybackslash}Xrrr@{}}
\toprule
Anomaly & How detected & How handled & TS & VS & \textbf{VS train}\\
\midrule
Image not stored axially (PIR, PSR, LAS) & orientation codes of the affine & resampled to RAS & 0 & 48 & \textbf{41}\\
Thick slices (more than 3$\times$ the in-plane size)$^{a}$ & voxel sizes & smoothed before downsampling & 0 & 4 & \textbf{0}\\
Only qform or only sform set & header codes & the one that is set is used & 1939 & 80 & \textbf{80}\\
CT stored as float, HU not whole numbers & data type and values & accepted (resampled by provider) & 0 & 77 & \textbf{73}\\
Label and CT affines differ by $\le$\,0.001\,mm & affine difference & treated as the same grid & 228 & 36 & \textbf{28}\\
Label mask stored as float & data type & rounded to whole numbers & 0 & 1 & \textbf{0}\\
\midrule
Two masks claim the same voxel (mostly S1 and L5) & overlap of the masks & the lumbar level wins & 817 & 0 & \textbf{0}\\
Padding below $-2000$\,HU outside the scan circle & share of such voxels & set to air by the HU window & 104 & 0 & \textbf{0}\\
Metal near the spine & $>$\,0.5\,ml above 3000\,HU & clipped at 1500\,HU, flagged & 124 & 4 & \textbf{4}\\
Vertebra cut off by the scan edge & mask touches the border & kept; Dice also reported without & 431 & 8 & \textbf{4}\\
No lumbar vertebra in the scan & label content & excluded from training & 720 & 13 & \textbf{15}\\
Sixth lumbar vertebra (L6) & VerSe label 25 & own class; excluded from training$^{b}$ & 0 & 3 & \textbf{3}\\
\bottomrule
\end{tabularx}
\end{table}

\begin{table}[h]
\centering\footnotesize
\caption{Dice per vertebra on the two test sets, mean / median (Table~2 of the original report, with a row for each model). \emph{Limit}: Dice of the ground truth after resampling to the model's voxel size and back, i.e.\ the highest score possible at that voxel size. $n$: number of scored vertebrae.}
\label{tab:dice}
\begin{tabular}{@{}llcccccccc@{}}
\toprule
Test set & Trained on (voxel size) & L1 & L2 & L3 & L4 & L5 & All & Limit & $n$\\
\midrule
TotalSegmentator test (78 scans) & TotalSegmentator (902 scans) (1.5\,mm) & .88 / .97 & .90 / .97 & .92 / .97 & .91 / .97 & .87 / .96 & \textbf{.89 / .97} & 1.00 & 328\\
 & VerSe19 training set (56 scans) (1\,mm) & \multicolumn{8}{c}{\emph{not evaluated yet}}\\
\midrule
VerSe19 val + test (67 scans) & TotalSegmentator (902 scans) (1.5\,mm) & .90 / .94 & .91 / .94 & .92 / .94 & .91 / .94 & .89 / .93 & \textbf{.91 / .94} & .96 & 305\\
 & VerSe19 training set (56 scans) (1\,mm) & .58 / .77 & .72 / .94 & .83 / .94 & .83 / .94 & .86 / .93 & \textbf{.76 / .93} & 1.00 & 305\\
\bottomrule
\end{tabular}
\end{table}

\begin{table}[h]
\centering\footnotesize
\caption{Further comparison numbers. \emph{Counted one level off}: renaming all predicted levels by one raises the scan's mean Dice by more than 0.1. \emph{Whole-vertebra false positives}: predicted levels absent from the ground truth and at least 20\,ml large. \emph{Body centre}: distance between the predicted and the ground-truth vertebral body centre, median / 90th percentile.}
\label{tab:extra}
\begin{tabular}{@{}lcccc@{}}
\toprule
 & \multicolumn{2}{c}{TotalSegmentator test (78 scans)} & \multicolumn{2}{c}{VerSe19 val + test (67 scans)}\\
Trained on & TotalSegmentator & VerSe19 train & TotalSegmentator & VerSe19 train\\
\midrule
Dice, complete vertebrae only (mean) & .92 & -- & .91 & .77\\
Scans counted one level off & 2 & -- & 1 & 7\\
Scans with a false-positive level & 52 & -- & 23 & 23\\
Whole-vertebra false positives & 11 & -- & 26 & 72\\
Body centre distance (mm) & 0.3 / 15.7 & -- & 0.6 / 3.9 & 0.9 / 34.3\\
\bottomrule
\end{tabular}
\end{table}


\section{Findings}
\begin{itemize}
\item \textbf{On a typical VerSe scan both models segment equally well} (median Dice .93 vs .94), although
  the VerSe-trained model saw only 56 training scans; at 1\,mm VerSe also loses nothing to resampling
  (limit 1.00 vs .96).
\item \textbf{On average the VerSe-trained model is clearly worse} (mean Dice .76 vs .91). The loss is in
  \emph{naming}, not in outlining: 7 scans are counted one level off (TotalSegmentator-trained: 1), and L1
  suffers most (mean .58 vs .90).
\item \textbf{It labels neck and upper thoracic vertebrae as lumbar far more often}: 72 whole-vertebra false
  positives against 26. Underneath them in the ground truth lie neck vertebrae C1--C7 in 32 cases
  (TotalSegmentator-trained: 8), thoracic vertebrae in 13 (5), unlabelled bone in 21 (12) and real lumbar
  vertebrae with the wrong level name in 6 (1). Many VerSe test scans show the whole spine. Both models were
  trained on lumbar crops only (T12 to sacrum), but with 56 instead of 902 scans the VerSe-trained model has
  learned far less well what is \emph{not} a lumbar vertebra. As the clean-up keeps only the largest piece per
  class, a wrongly named neck vertebra can even replace the real L1.
\item \textbf{Body centres}: the median distance is similar (0.9 vs 0.6\,mm), but the 90th percentile is
  34\,mm against 4\,mm: the same naming errors put a ``body centre'' on the wrong vertebra.
\item \textbf{Takeaway}: here, many training scans from another source beat few scans from the target source.
  56 scans are not enough to learn level counting on whole-spine scans; more scans, training patches outside
  the lumbar crop, or a rule-based level check would be the next steps.
\end{itemize}
\textit{The TotalSegmentator test set is not yet evaluated for the VerSe-trained model (60 of 109 scans done);
its rows are filled in automatically once the evaluation is complete.}

\end{document}
